\documentclass[10pt]{article}

\usepackage[utf8]{inputenc}

\usepackage[T1]{fontenc}

\usepackage{amsmath}

\usepackage{amsfonts}

\usepackage{amssymb}

\usepackage[version=4]{mhchem}

\usepackage{stmaryrd}



\begin{document}

$y, z)$ с помощью о г р а ни ч ен н о го к в а т т о р а, если



$P\left(x_{1}, \ldots, x_{n}, z\right) \Leftrightarrow \forall y\left(y \leqslant z \Rightarrow Q\left(x_{1}, \ldots, x_{n}, y, z\right)\right)$



или



$P\left(x_{1}, \ldots, x_{n}, z\right) \Leftrightarrow \exists y\left(y \leqslant z \& Q\left(x_{1}, \ldots, x_{n}, y, z\right)\right)$.



Класс п. р. о. замкнут относительно применения всех логич. связок (включая отрицание) и ограниченных кванторов.



Пусть $f_{1}, \ldots, f_{s+1}$ суть $n$-местные П. р. ф., а $P_{1}, \ldots$ $\ldots, P_{s}$ - такие п. p. о., что на любом наборе значений аргументов истинно не более одного из них. Тогда функция



$f\left(x_{1}, \ldots, x_{n}\right)=\left\{\begin{array}{c}f_{1}\left(x_{1}, \ldots, x_{n}\right), \text { если } P_{1}\left(x_{1}, \ldots, x_{n}\right), \\ \cdot . \cdot . \cdot . \cdot . \cdot . \cdot . \cdot . \\ f_{s}\left(x_{1}, \ldots, x_{n}\right), \text { если } P_{s}\left(x_{1}, \ldots, x_{n}\right), \\ f_{s+1}\left(x_{1}, \ldots, x_{n}\right), \text { в других случаях, }\end{array}\right.$ является П. р.ф.



Говорят, что функция $f\left(x_{1}, \ldots, x_{n}, z\right)$ получена из всюду определенной функции $g\left(x_{1}, \ldots, x_{n}, y, z\right)$ примевением о г р а и и е н но г о о п е р а то р а м ив и м и з а ц и и, если $f\left(x_{1}, \ldots, x_{n}, z\right)$ равно минимальному $y$ такому, что $y \leqslant z$ и $g\left(x_{1}, \ldots, x_{n}, y, z\right)=0$, и равно $z+1$, если такого $y$ нет. Класс П. р. Ф. замкнут относительно применения ограниченных операторов минимизации.



Функция Ф $\left(y, x_{1}, \ldots, x_{n}\right)$ наз. у н и в е р с а л ь н о й для класса всех $n$-местных П. р. ф., если для каждой П. р. Ф. $f\left(x_{1}, \ldots, x_{n}\right)$ найдется натуральное число $k$ такое, что



\[

f\left(x_{1}, \ldots, x_{n}\right)=\Phi\left(k, x_{1}, \ldots, x_{n}\right) .

\]



Для каждого $n \geqslant 1$ такая универсальная функция существует, но она не может быть П. р. ф.



Всякое рекурсивно перечислимое множество есть область значений П. р. Ф.; всякоө рекурсивно перечислимое отношение $R\left(x_{1}, \ldots, x_{n}\right)$ представимо в виде $\exists y A\left(y, x_{1}, \ldots, x_{n}\right)$, где $A$ - п. р. о. Всякая П. р. ф. представима в арифметике формальной, т. е. для каждой П. р. ф. $f\left(x_{1}, \ldots, x_{n}\right)$ найдется арифметич. формула $F\left(y, x_{1}, \ldots, x_{n}\right)$ такая, что для натуральных $k_{1}, \ldots, k_{n}$, $m$ при $f\left(k_{1}, \ldots, k_{n}\right)=m$ в формальной арифметике выводима формула $F\left(\bar{m}, \bar{k}_{1}, \ldots, \bar{k}_{n}\right)$, а при $f\left(k_{\underline{1}}, \ldots, k_{n}\right) \neq m$ выводима $\rceil F\left(\bar{m}, \bar{k}_{1}, \ldots, \bar{k}_{n}\right)$ (здесь $\bar{k}_{1}, \ldots, \bar{k}_{n}, \bar{m}-$ арифметич. термы, изображаюпце в формальной арифметике натуральные числа $\left.k_{1}, \ldots, k_{n}, m\right)$. Этот факт занимает центральное место в доказательстве неполноты формальной арифметики (см. [4]).



Лит.: [1] Ус п е н с к и й В. А., Лекции о вычислимых функциях, М., 1960; [2] М а льц е в А. И., Алгоритмы и рекурсивные функции, М., 1965; [3] Р о д ж е р с Х., Теория рекурсивных функций п эффективная вычислимость, пер. с англ., М., 1972; [4̂] М е н д е л ь с о н Э., Введение в математическую логику, пер. с англ., М., 1976.



ПРИМИТИВНОЕ КОЛЬЦО п р а в о е - ассоциативное кольцо, обладающее правым точным неприводимы.м модулем. Аналогично (с помощью левого неприводимого модуля) определяется л е в о е п р и м и т и вн о к о л ц о. Классы правых и левых П. к. не совпадают. Всякое коммутативное П. к. является полем. Всякое полупростое (в смысле Джекобсона радикала) кольцо является подпрямым произведеныем П. к. Простое кольцо либо является П. к., либо радикально. П. к. с ненулевыми минимальными правыми идеалами описываются теоремой плотности. П. к. с условием минимальности для правых идеалов (т. е. артиновы П. к.) являются простыми .



Кольцо $R$ примитивно тогда и только тогда, когда оно обладает максимальным модулярным правым идеалом $I$, к-рый не содержит двусторонних идеалов кольца $R$, отличных от нулевого идеала. Әто свойство может быть принято за определение П. к. в классе неассоциативных колец.



Лит.: [1] Д ж е к о б с о н Н., Строеве колец, пер. с англ., М., 1961; [2] Х е р с т е й в И., Некоммутативные кольца, пер. ПРИМИТИВНЫЫ ИДЕАЛ, п р а в о п р и м и т и вны й и д е а л, - такой двусторонний идеал $P$ ассоциативного кольца $R$, что факторкольцо $R / P$ является (правым) примитивныл кольцом. Аналогично, с помощью левых примитивных колец может быть определен л е в о п р и м и т и в ны й и д е а л. Множество П всех П. и. кольца, снабженное нек-рой топологией, оказывается полезным при изучении отдельных классов колец. Обычно множество П топологизируется при помощи следующего замыкания отношения:



\[

A=\left\{P^{\prime} \mid P^{\prime} \in \Pi, P^{\prime} \supseteq(\cap P, P \in A)\right\}

\]



где $A$ подмножество в П. Множество всех П. и. кольца, снабженное такой топологией, наз. с т р у к т у р ны м п р о с т р а н с т в о м әтого кольца.



Лит.: [1] Д ж е к о б с о н Н., Строение колец, пер. с англ., М., 1961.



ПРИМИТИВНЫИ КЛАСС а л г е б р а ич е с к й с и с т е м - то же, что многообразие (см. Алгебраических систем многообразие).



ПРИМИТИВНЫЙ МНОГОчЛЕН - многочлен $f(X) \in R[X]$, где $R-$ ассоциативно-коммутативное кольцо с оцнозначным разложением на множители, коэффициенты к-рого не имеют нетривиальных общих делителей. Любой многочлен $g(X) \in R[X]$ можно записать в виде $g(X)=c(g) f(X)$, где $f(X)-\Pi$. м., а $c(g)$ наибольший общий делитель коәффициентов многочлена $g(X)$. Элемент $c(g) \in R$, определенный с точностью до умножения на обратимые элементы из $R$, наз. с од е р жа ние м мно о очлен а $g(X)$. Справедлива л е м м а $\Gamma$ а у с с а: если $g_{1}(X), g_{2}(X) \in R[X]$, то $c\left(g_{1} g_{2}\right)=c\left(g_{1}\right) c\left(g_{2}\right)$. В частности, произведение П. М. снова примитивно.



Лит.: [1] 3 а р и с с к и й О., С а м ю ә л ь П., Коммутативная алгебра, пер. с англ., т. 1, М., 1963. Л. В. Кузьмин.



ПРИОРИТЕТА МЕТОД - метод, применяемый в рекурсивной теории множеств для построения просто устроенных с рекурсивной точки зрения (в простейших случаях - рекурсивно перечислимых) множеств (функций, нумераций и т. П.), удовлетворяющих бесконечной системе условий определенного типа. Специфика допустимых условий такова, что для того чтобы удовлетворить отдельному условию из данной системы, обычно бывает достаточно, чтобы в строящеөся множество был включен определенный (зависящий от рассматриваемого условия) әлемент. Однако на каждом этапе построения (представляющего собой нек-рый вычислительный процесс, чөм и обеспөчивается рекурсивная простота строения искомого объекта) каждое из условий системы (к-рое определяется, вообще говоря, бесконечным множеством конструктивных объектов) представлено нек-рой своей конөчной аппроксимацией. Включение в строящееся множество элемента, обеспечивающего выполнение аппроксимации $j$-го условия, ещце не дает гарантии, что тем самым будет удовлетворено само $j$-е условие. П. м. позволяет в известных случаях обходить это препятствие. С этой целью j-му условию данной системы, $j=1,2,3, \ldots$, в процессе построения ставится в соответствие натуральное число, являющееся кандидатом на роль элемента, включение к-рого в строяшееся множество удовлетворяет $j$-му условию; про такой элемент принято говорить, что он помечен маркером $|\bar{j}|$ (снабженным, возможно, дополнительными индексами). Каждый такой кандидат может в процессе построения заменяться (или, что то же самое, маркер мокет сдвигаться), но при этом в простейших приоритетных конструкциях сама после-





\end{document}